% Part: incompleteness
% Chapter: arithmetization-syntax
% Section: coding-terms

\documentclass[../../../include/open-logic-section]{subfiles}

\begin{document}

% SEGMENT OLP-0283-S01

\olfileid{inc}{art}{trm}
\olsection{சொற்களைக் குறியீடாக்குதல்}

\begin{explain}
ஒரு சொல் என்பது குறிப்பிட்ட வகையைச் சேர்ந்த குறிகளின் ஒரு தொடரே: சொற்களுக்கான
அமைப்புவிதிகளின்படி மாறிலிகளிலிருந்தும் மாறிகளிலிருந்தும் அது
தொகுத்தறிதல் முறையில் கட்டப்படுகிறது. குறிகளின் தொடர்களை எண்களாகக்
குறியீடாக்கலாம் என்பதால்---குறிகளுக்கான ஒரு குறியீட்டு முறையையும் எண்களின்
தொடர்களைக் குறியீடாக்கும் ஒரு வழியையும் பயன்படுத்தி---சொற்களுக்குக் கோடல்
எண்களை ஒதுக்குவது கடினமல்ல. மாறாக, செம்மையாக அமைந்த ஒரு சொல்லின் கோடல்
எண்ணாக ஓர் எண் இருப்பது என்ற பண்பு கணிக்கத்தக்கது, உண்மையில் முதன்மை
மீள்வரையறையானது, என்று காட்டுவதே சவால்.
\end{explain}

!!^{variable}s மற்றும் !!{constant}s ஆகியவையே மிக எளிய சொற்கள்; $x$ ஆனது
அத்தகைய ஒரு சொல்லின் கோடல் எண்ணா என்று சோதிப்பது எளிது: ஏதாவது~$i$ க்கு
$x$ ஆனது~$\Gn{\Obj v_i}$ ஆக இருந்தால், $\fn{Var}(x)$ பொருந்தும். வேறு
சொற்களில், $x$~ஆனது நீளம்~$1$ கொண்ட ஒரு தொடர்; அதன் ஒரே உறுப்பு~$(x)_0$
ஏதாவது !!{variable}~$\Obj v_i$ இன் குறியீடு. அதாவது, ஏதாவது~$i$ க்கு
$x$ ஆனது~$\tuple{\tuple{1, i}}$. இதேபோல, ஏதாவது~$i$ க்கு $x$ ஆனது
$\Gn{\Obj c_i}$ ஆக இருந்தால், $\fn{Const}(x)$ பொருந்தும். இவ்விரு
தொடர்புகளும் முதன்மை மீள்வரையறையானவை; ஏனெனில் அத்தகைய~$i$ ஒன்று இருந்தால்,
அது $< x$ ஆக இருக்க வேண்டும்:
\begin{align*}
  \fn{Var}(x) & \defiff \bexists{i<x}{x = \tuple{\tuple{1, i}}}\\
  \fn{Const}(x) & \defiff \bexists{i<x}{x = \tuple{\tuple{2, i}}}
\end{align*}

\begin{prop}
\ollabel{prop:term-primrec}
$x$ ஆனது முறையே ஒரு சொல்லின் அல்லது ஒரு மூடிய சொல்லின் கோடல் எண்ணாக
இருந்தால், அப்போதும் அப்போது மட்டுமே பொருந்தும் $\fn{Term}(x)$ மற்றும்
$\fn{ClTerm}(x)$ என்ற தொடர்புகள் முதன்மை மீள்வரையறையானவை.
\end{prop}

\begin{proof}
குறிகளின் ஒரு தொடர்~$s$ ஒரு சொல்லாக இருந்தால், அப்போதும் அப்போது மட்டுமே
$s_0$, \dots, $s_{k-1} = s$ என்ற சொற்களின் தொடர் ஒன்று இருக்கும்; சொற்களின்
அமைப்புவிதிகளின்படி !!{constant}s மற்றும் !!{variable}s இலிருந்து~$s$
எவ்வாறு அமைக்கப்பட்டது என்பதை அத்தொடர் பதிவுசெய்யும். முன்வைக்கப்பட்ட
அத்தகைய ஓர் அமைப்புத் தொடர் அமைப்புவிதிகளைப் பின்பற்றுகிறது என்று கூற,
ஒவ்வொரு $i < k$ க்கும் பின்வருவனவற்றுள் ஒன்று பொருந்த வேண்டும்:
\begin{enumerate}
\item $s_i$ என்பது !!a{variable}~$\Obj v_j$, அல்லது
\item $s_i$ என்பது !!a{constant}~$\Obj c_j$, அல்லது
\item $i$-ஆம் இடத்துக்கு முன்பு வரும் $n$ சொற்களான $t_1$, \dots, $t_n$
  ஆகியவற்றிலிருந்து, ஓர் $n$-இட !!{function}~$\Obj f^n_j$ ஐப் பயன்படுத்தி
  $s_i$ கட்டப்படுகிறது.
\end{enumerate}
கோடல் எண்கள் மீது இதற்கு ஒத்த தொடர்பு முதன்மை மீள்வரையறையானது என்று காட்ட,
இந்நிபந்தனையை முதன்மை மீள்வரையறை முறையில், அதாவது முதன்மை மீள்வரையறைச்
சார்புகள், தொடர்புகள் மற்றும் வரம்பிட்ட அளவையிடல் ஆகியவற்றைப் பயன்படுத்தி,
வெளிப்படுத்த வேண்டும்.

$s_0$, \dots, $s_{k-1}$ என்ற தொடரைக் குறியீடாக்கும் எண்~$y$ என்று கொள்க;
அதாவது, $y = \tuple{\Gn{s_0}, \dots, \Gn{s_{k-1}}}$. ஒவ்வொரு $i < k$
க்கும் பின்வருவன பொருந்தினால், அப்போதும் அப்போது மட்டுமே, கோடல் எண்~$x$
கொண்ட சொல்லுக்கான ஓர் அமைப்புத் தொடரை அது குறியீடாக்குகிறது:
\begin{enumerate}
\item $\fn{Var}((y)_i)$, அல்லது
\item $\fn{Const}((y)_i)$, அல்லது
\item ஓர்~$n$ உம் $z = \tuple{z_1, \dots, z_n}$ என்ற எண்ணும் உள்ளன;
  ஒவ்வொரு~$z_l$ உம் ஏதாவது $i' < i$ க்கு $(y)_{i'}$ உடன் சமம், மேலும்
\[
(y)_i = \Gn{\Obj f^n_j(} \concat \fn{flatten}(z) \concat \Gn{)},
\]
\end{enumerate}
மேலும் $(y)_{k-1} = x$. ($\fn{flatten}(z)$ என்ற சார்பு,
$\tuple{\Gn{t_1}, \dots, \Gn{t_n}}$ என்ற தொடரை $\Gn{t_1, \dots,
  t_n}$ ஆக மாற்றுகிறது; அது முதன்மை மீள்வரையறையானது.)

(3) இல் உள்ள சுட்டெண்கள் $j$, $n$, சொற்கள்~$t_l$ இன் கோடல் எண்கள்~$z_l$,
மற்றும் $\tuple{z_1, \dots, z_n}$ என்ற தொடரின் குறியீடு~$z$ ஆகிய அனைத்தும்
~$y$ ஐவிடச் சிறியவை. மேலே உள்ள $k$ ஐ $\len{y}$ கொண்டு பதிலிடலாம். ஆகவே,
``கோடல் எண்~$x$ கொண்ட சொல்லின் ஓர் அமைப்புத் தொடரின் குறியீடே~$y$'' என்பதை
இந்தத் தொடர்பு முதன்மை மீள்வரையறையானது என்று காட்டும் வகையில் வெளிப்படுத்தலாம்.

இப்போது~$y$ மீது ஒரு முதன்மை மீள்வரையறை வரம்பு உள்ளது என்பதை மட்டும் நாம்
உறுதிப்படுத்த வேண்டும். ஆனால் $x$ ஒரு சொல்லின் கோடல் எண்ணாக இருந்தால், அது
அதிகபட்சம் $\len{x}$ சொற்களைக் கொண்ட ஓர் அமைப்புத் தொடரைக் கொண்டிருக்க
வேண்டும். (ஏனெனில்~$s$ இன் அமைப்புத் தொடரிலுள்ள ஒவ்வொரு சொல்லும்~$s$ இல்
ஏதாவது ஓர் இடத்தில் தொடங்க வேண்டும்; எந்த இரண்டு உட்சொற்களும் ஒரே இடத்தில்
தொடங்க முடியாது.) $s$ இன் ஒவ்வொரு உட்சொல்லின் கோடல் எண்ணும் நிச்சயமாக
$\le x$. எனவே, $k=\len{x}$ எனும்போது, $\le p_{k-1}^{k(x+1)}$ என்ற
குறியீட்டைக் கொண்ட ஓர் அமைப்புத் தொடர் எப்போதும் உள்ளது.

$\fn{ClTerm}$ க்கு, !!^{variable}s பற்றிய கூறை மட்டும் நீக்கிவிடுக.
\end{proof}

\begin{prob}
$\tuple{\Gn{t_1}, \dots, \Gn{t_n}}$ என்ற தொடரை $\Gn{t_1, \dots, t_n}$
ஆக மாற்றும் $\fn{flatten}(z)$ என்ற சார்பு முதன்மை மீள்வரையறையானது என்று
காட்டுக.
\end{prob}

\begin{prop}\ollabel{prop:num-primrec}
  $\fn{num}(n) = \Gn{\num{n}}$ என்ற சார்பு முதன்மை மீள்வரையறையானது.
\end{prop}

\begin{proof}
  $\fn{num}(n)$ ஐ முதன்மை மீள்வரையறையால் வரையறுக்கிறோம்:
  \begin{align*}
    \fn{num}(0) & = \Gn{\Obj 0}\\
    \fn{num}(n+1) & = \Gn{\prime(} \concat \fn{num}(n) \concat \Gn{)}.
  \end{align*}
\end{proof}

\end{document}
